\slajdid{02-s043}
\begin{frame}[label=02-s043]
\gusttekst
Primer: Neka je na osnovu zahteva za kombinacionom mrežom napravljena Karnoova tabela

\begin{columns}[c,onlytextwidth]
\begin{column}{.38\textwidth}\centering
% Bitovi indeksa su DCBA; kolone BA=00,01,11,10, redovi DC=00,01,11,10.
\begingroup
\fontsize{16.67}{18.5}\selectfont
\renewcommand{\small}{\fontsize{16.67}{18.5}\selectfont}
\scalebox{.54}{%
\begin{karnaugh-map}(label=middle)[4][4][1][$A$][$B$][$C$][$D$]
\minterms{0,4,5,8,10}
\terms{1}{b}
\maxterms{2,3,6,7,9,11,12,13,14,15}
\end{karnaugh-map}
}
\endgroup

\end{column}
\begin{column}{.59\textwidth}\gusttekst
Da bi dobili minimalnu funkciju u obliku zbira proizvoda prekrivamo sve logičke jedinice površinama najvišeg reda, koristeći se i sa stanjem b
\end{column}
\end{columns}
\begin{columns}[c,onlytextwidth]
\begin{column}{.38\textwidth}\centering
\begingroup
\fontsize{16.67}{18.5}\selectfont
\renewcommand{\small}{\fontsize{16.67}{18.5}\selectfont}
\scalebox{.54}{%
\begin{karnaugh-map}(label=middle,implicantcolors={DEplava,DEcrvena})[4][4][1][$A$][$B$][$C$][$D$]
\minterms{0,4,5,8,10}
\terms{1}{b}
\maxterms{2,3,6,7,9,11,12,13,14,15}
\implicant{0}{5}
\implicantedge{8}{8}{10}{10}
\end{karnaugh-map}
}
\endgroup

\end{column}
\begin{column}{.59\textwidth}\gusttekst
Pošto suštinski površine predstavljaju objedinjavanje proizvoda kod kojih je moguće uraditi sažimanje, u proizvodima ostaju literali koji se \alert{\underline{ne menjaju}} na površini.
\[
F=\bar{D}\bar{B}+D\bar{C}\bar{A}
\]
\end{column}
\end{columns}
\end{frame}
